\documentclass{amsart}
% For dealing with references we use the comment environment
\usepackage{verbatim}
\newenvironment{reference}{\comment}{\endcomment}
%\newenvironment{reference}{}{}
\newenvironment{slogan}{\comment}{\endcomment}
\newenvironment{history}{\comment}{\endcomment}

% For commutative diagrams we use Xy-pic
\usepackage[all]{xy}

% We use 2cell for 2-commutative diagrams.
\xyoption{2cell}
\UseAllTwocells

% We use multicol for the list of chapters between chapters
\usepackage{multicol}

% This is generally recommended for better output
\usepackage{lmodern}
\usepackage[T1]{fontenc}

% For cross-file-references

% Package for hypertext links:
\usepackage{hyperref}

% For any local file, say "hello.tex" you want to link to please

% Theorem environments.
%
\theoremstyle{plain}
\newtheorem{theorem}[subsection]{Theorem}
\newtheorem{proposition}[subsection]{Proposition}
\newtheorem{lemma}[subsection]{Lemma}

\theoremstyle{definition}
\newtheorem{definition}[subsection]{Definition}
\newtheorem{example}[subsection]{Example}
\newtheorem{exercise}[subsection]{Exercise}
\newtheorem{situation}[subsection]{Situation}

\theoremstyle{remark}
\newtheorem{remark}[subsection]{Remark}
\newtheorem{remarks}[subsection]{Remarks}

\numberwithin{equation}{subsection}

% Macros
%
\def\lim{\mathop{\mathrm{lim}}\nolimits}
\def\colim{\mathop{\mathrm{colim}}\nolimits}
\def\Spec{\mathop{\mathrm{Spec}}}
\def\Hom{\mathop{\mathrm{Hom}}\nolimits}
\def\Ext{\mathop{\mathrm{Ext}}\nolimits}
\def\SheafHom{\mathop{\mathcal{H}\!\mathit{om}}\nolimits}
\def\SheafExt{\mathop{\mathcal{E}\!\mathit{xt}}\nolimits}
\def\Sch{\mathit{Sch}}
\def\Mor{\mathop{\mathrm{Mor}}\nolimits}
\def\Ob{\mathop{\mathrm{Ob}}\nolimits}
\def\Sh{\mathop{\mathit{Sh}}\nolimits}
\def\NL{\mathop{N\!L}\nolimits}
\def\CH{\mathop{\mathrm{CH}}\nolimits}
\def\proetale{{pro\text{-}\acute{e}tale}}
\def\etale{{\acute{e}tale}}
\def\QCoh{\mathit{QCoh}}
\def\Ker{\mathop{\mathrm{Ker}}}
\def\Im{\mathop{\mathrm{Im}}}
\def\Coker{\mathop{\mathrm{Coker}}}
\def\Coim{\mathop{\mathrm{Coim}}}

% Boxtimes
%
\DeclareMathSymbol{\boxtimes}{\mathbin}{AMSa}{"02}

%
% Macros for moduli stacks/spaces
%
\def\QCohstack{\mathcal{QC}\!\mathit{oh}}
\def\Cohstack{\mathcal{C}\!\mathit{oh}}
\def\Spacesstack{\mathcal{S}\!\mathit{paces}}
\def\Quotfunctor{\mathrm{Quot}}
\def\Hilbfunctor{\mathrm{Hilb}}
\def\Curvesstack{\mathcal{C}\!\mathit{urves}}
\def\Polarizedstack{\mathcal{P}\!\mathit{olarized}}
\def\Complexesstack{\mathcal{C}\!\mathit{omplexes}}
% \Pic is the operator that assigns to X its picard group, usage \Pic(X)
% \Picardstack_{X/B} denotes the Picard stack of X over B
% \Picardfunctor_{X/B} denotes the Picard functor of X over B
\def\Pic{\mathop{\mathrm{Pic}}\nolimits}
\def\Picardstack{\mathcal{P}\!\mathit{ic}}
\def\Picardfunctor{\mathrm{Pic}}
\def\Deformationcategory{\mathcal{D}\!\mathit{ef}}

\setlength{\emergencystretch}{3em}
\begin{document}
\title[Stacks Project, Tag 07FG]{733 --- Artinian local algebras in characteristic p with finite cotangent H-one are unions of flat finite-type models}
\maketitle
\noindent Copyright (C) 2005--2025 Johan de Jong. Stacks Project authors. Source: \href{https://stacks.math.columbia.edu/tag/07FG}{Tag 07FG}.

\noindent Editorial difficulty note (not author text). Difficulty H3: The proof adjusts coefficient-field sections by derivations, uses finite-dimensional cotangent homology to confine the obstruction to finitely many differential coordinates, and inducts through nilpotent layers. This specialized inseparability and approximation construction is research-level commutative algebra..

\noindent Editorial provenance note (not author text). History: excerpt from revision \texttt{a0444\allowbreak 6e57e\allowbreak c1fbc\allowbreak 252a8\allowbreak 71afc\allowbreak ec775\allowbreak 2fb28\allowbreak 07b14}, smoothing.tex, lines 2444--2549. Reference presentation was adapted; original-author context and core support blocks were retained or added. The mathematical wording is unchanged. Licensed under GNU FDL 1.2 or later; no invariant sections or cover texts. See ../licenses/COPYING. Context blocks reproduce author definitions and supporting results; they are not additional counted problems.

\begin{lemma}
\label{lemma-helper}
Let $k$ be a field of characteristic $p > 0$.
Let $(\Lambda, \mathfrak m, K)$ be an Artinian local $k$-algebra.
Assume that $\dim H_1(L_{K/k}) < \infty$.
Then $\Lambda$ is a filtered colimit of Artinian
local $k$-algebras $A$ with each map $A \to \Lambda$ flat, with
$\mathfrak m_A \Lambda = \mathfrak m$, and with
$A$ essentially of finite type over $k$.
\end{lemma}

\begin{proof}
Note that the flatness of $A \to \Lambda$ implies that $A \to \Lambda$
is injective, so the lemma really tells us that $\Lambda$ is a
directed union of these types of subrings $A \subset \Lambda$.
Let $n$ be the minimal integer such that $\mathfrak m^n = 0$.
We will prove this lemma by induction on $n$. The case $n = 1$ is clear
as a field extension is a union of finitely generated field extensions.

\medskip\noindent
Pick $\lambda_1, \ldots, \lambda_d \in \mathfrak m$ which generate
$\mathfrak m$. As $K$ is formally smooth over $\mathbf{F}_p$ (see
Algebra, Lemma \href{https://stacks.math.columbia.edu/tag/0320}{lemma formally smooth extensions easy (Tag 0320)}) we can
find a ring map $\sigma : K \to \Lambda$ which is a section of the
quotient map $\Lambda \to K$. In general $\sigma$ is {\bf not}
a $k$-algebra map. Given $\sigma$ we define
$$
\Psi_\sigma : K[x_1, \ldots, x_d] \longrightarrow \Lambda
$$
using $\sigma$ on elements of $K$ and mapping $x_i$ to $\lambda_i$.
Claim: there exists a $\sigma : K \to \Lambda$
and a subfield $k \subset F \subset K$ finitely generated over $k$
such that the image of $k$ in $\Lambda$ is contained in
$\Psi_\sigma(F[x_1, \ldots, x_d])$.

\medskip\noindent
We will prove the claim by induction on the least integer $n$ such that
$\mathfrak m^n = 0$. It is clear for $n = 1$. If $n > 1$ set
$I = \mathfrak m^{n - 1}$ and $\Lambda' = \Lambda/I$.
By induction we may assume
given $\sigma' : K \to \Lambda'$ and $k \subset F' \subset K$ finitely
generated such that the image of $k \to \Lambda \to \Lambda'$
is contained in $A' = \Psi_{\sigma'}(F'[x_1, \ldots, x_d])$.
Denote $\tau' : k \to A'$ the induced map.
Choose a lift $\sigma : K \to \Lambda$ of $\sigma'$ (this is possible
by the formal smoothness of $K/\mathbf{F}_p$ we mentioned above).
For later reference we note that we can change $\sigma$ to
$\sigma + D$ for some derivation $D : K \to I$.
Set $A = F[x_1, \ldots, x_d]/(x_1, \ldots, x_d)^n$.
Then $\Psi_\sigma$ induces a ring map
$\Psi_\sigma : A \to \Lambda$. The composition with the
quotient map $\Lambda \to \Lambda'$ induces a surjective
map $A \to A'$ with nilpotent kernel.
Choose a lift $\tau : k \to A$ of $\tau'$ (possible as $k/\mathbf{F}_p$
is formally smooth). Thus we obtain two maps $k \to \Lambda$, namely
$\Psi_\sigma \circ \tau : k \to \Lambda$ and the given map $i : k \to \Lambda$.
These maps agree modulo $I$, whence the difference is a
derivation $\theta = i - \Psi_\sigma \circ \tau : k \to I$.
Note that if we change $\sigma$ into $\sigma + D$ then we change
$\theta$ into $\theta - D|_k$.

\medskip\noindent
Choose a set of elements $\{y_j\}_{j \in J}$ of $k$ whose differentials
$\text{d}y_j$ form a basis of $\Omega_{k/\mathbf{F}_p}$. The Jacobi-Zariski
sequence for $\mathbf{F}_p \subset k \subset K$ is
$$
0 \to H_1(L_{K/k}) \to \Omega_{k/\mathbf{F}_p} \otimes K \to
\Omega_{K/\mathbf{F}_p} \to \Omega_{K/k} \to 0
$$
As $\dim H_1(L_{K/k}) < \infty$ we can find a finite subset $J_0 \subset J$
such that the image of the first map is contained in
$\bigoplus_{j \in J_0} K\text{d}y_j$. Hence the elements
$\text{d}y_j$, $j \in J \setminus J_0$ map to $K$-linearly independent
elements of $\Omega_{K/\mathbf{F}_p}$. Therefore we can choose
a $D : K \to I$ such that $\theta - D|_k = \xi \circ \text{d}$
where $\xi$ is a composition
$$
\Omega_{k/\mathbf{F}_p} = \bigoplus\nolimits_{j \in J} k \text{d}y_j
\longrightarrow \bigoplus\nolimits_{j \in J_0} k \text{d}y_j
\longrightarrow I
$$
Let $f_j = \xi(\text{d}y_j) \in I$ for $j \in J_0$.
Change $\sigma$ into $\sigma + D$ as above. Then we see that
$\theta(a) = \sum_{j \in J_0} a_j f_j$ for $a \in k$ where
$\text{d}a = \sum a_j \text{d}y_j$ in $\Omega_{k/\mathbf{F}_p}$.
Note that $I$ is generated by the monomials
$\lambda^E = \lambda_1^{e_1} \ldots \lambda_d^{e_d}$ of
total degree $|E| = \sum e_i = n - 1$ in $\lambda_1, \ldots, \lambda_d$.
Write $f_j = \sum_E c_{j, E} \lambda^E$ with $c_{j, E} \in K$.
Replace $F'$ by $F = F'(c_{j, E})$. Then the claim holds.

\medskip\noindent
Choose $\sigma$ and $F$ as in the claim. The kernel of $\Psi_\sigma$ is
generated by finitely many polynomials
$g_1, \ldots, g_t \in K[x_1, \ldots, x_d]$ and we may assume their
coefficients are in $F$ after enlarging $F$ by adjoining finitely many
elements. In this case it is clear that the map
$A = F[x_1, \ldots, x_d]/(g_1, \ldots, g_t) \to
K[x_1, \ldots, x_d]/(g_1, \ldots, g_t) = \Lambda$ is flat.
By the claim $A$ is a $k$-subalgebra of $\Lambda$.
It is clear that $\Lambda$ is the filtered colimit of these
algebras, as $K$ is the filtered union of the subfields $F$.
Finally, these algebras are essentially of finite type over $k$ by
Algebra, Lemma
\href{https://stacks.math.columbia.edu/tag/07DT}{lemma essentially of finite type into artinian local (Tag 07DT)}.
\end{proof}
\end{document}
